\documentclass[12pt,a4paper]{article}
\usepackage[utf8]{inputenc}
\usepackage[T1]{fontenc}
\usepackage{amsmath,amssymb,amsthm}
\usepackage{graphicx}
\usepackage{geometry}
\usepackage{microtype}
\microtypesetup{expansion=false}
\usepackage{hyperref}
\usepackage{natbib}
\usepackage{booktabs}
\usepackage{titlesec}
\usepackage{float}
\usepackage{setspace}
\usepackage{parskip}
\usepackage{tikz}
\usepackage{pgfplots}
\pgfplotsset{compat=1.18}
\setcounter{tocdepth}{3}
\setcounter{secnumdepth}{3}

\geometry{margin=1in}

% Section formatting
\titleformat{\section}
  {\normalfont\Large\bfseries}
  {\thesection}
  {1em}
  {}

\titlespacing*{\section}
  {0pt}
  {12pt}
  {6pt}

% Subsection formatting
\titleformat{\subsection}
  {\normalfont\large\bfseries}
  {\thesubsection}
  {1em}
  {}

\titlespacing*{\subsection}
  {0pt}
  {8pt}
  {4pt}

\onehalfspacing

\title{Emotional Regulation as Control Signals: Active Homeostasis in Synthetic Cognitive Architectures}

\author{John H. Cragin \\
Independent Researcher \\
john.cragin@outlook.com}

\date{\today}

\begin{document}

\maketitle

Emotional states in biological cognition serve as control signals that modulate processing intensity, pathway activation, and behavioral adaptation. This paper demonstrates how SpiralBrain v3.0 implements emotions as active computational substrates, not peripheral classifiers, achieving anticipatory regulation with 89\% effectiveness. Using SEC vectors (valence, arousal, hazard pressure, neuromodulator flexibility) as foundational control parameters, the system maintains cognitive stability through predictive intervention and adaptive homeostasis. No claims of universality or biological equivalence are made; conclusions are constrained to empirically observed behavior in SpiralBrain.

This work builds upon the Regulatory Intelligence (RI) paradigm \cite{Cragin2026RI}, which defines intelligence as the capacity for viability rather than calculation. In the RI framework, emotions serve as the primary computational substrate shaping cognitive processing, with SEC vectors functioning as active regulatory mechanisms within a 128-dimensional cognitive manifold.

\section{Introduction}

In biological systems, emotions function as control signals that regulate cognitive processing, not merely as affective states or response classifiers. This regulatory role enables adaptive behavior, homeostasis preservation, and context-appropriate processing modulation. Traditional AI approaches treat emotions as optional features or post-processing classifiers, missing their fundamental role as computational substrates.

This paper examines how SpiralBrain v3.0 implements emotions as active control signals through SEC (Symbolic-Emotional Calibration) vectors, achieving anticipatory emotional regulation with measurable effectiveness. The system demonstrates a closed cognitive loop where emotional states predict, intervene, learn, and adapt—creating genuine synthetic emotional intelligence. All findings are grounded in executable Python artifacts with reproducible results.

The RI paradigm \cite{Cragin2026RI} provides the theoretical foundation: intelligence emerges from regulatory capability rather than computational scale, with emotions serving as the governing computational substrate that regulates, shapes, and stabilizes cognitive processing.

\section{The Regulatory Intelligence Framework}

SpiralBrain v3.0 implements the Regulatory Intelligence (RI) paradigm \cite{Cragin2026RI}, where intelligence is defined as the capacity for viability—the ability to maintain internal coherence and functional stability under stress. This paradigm inverts traditional AI design: instead of optimizing for external task accuracy, the system is mandated to maintain internal homeostasis, with specialized reasoning emerging as a strategy for stability preservation.

The architecture employs a four-lobe topology to prevent logical bleeding between regulatory and analytical functions:

\begin{itemize}
\item \textbf{Cortex (Metacognitive Lobe)}: Supervises temporal consistency and self-narrative
\item \textbf{Codex (Analytical Lobe)}: Processes symbolic logic and domain reasoning  
\item \textbf{Nexus (Affective Lobe)}: Generates SEC emotional signals as regulatory mechanisms
\item \textbf{Sensus (Perceptual Lobe)}: Interfaces with environmental data streams
\end{itemize}

Emotions function as the primary computational substrate within this framework, with SEC vectors serving as active regulatory parameters that modulate cognitive processing across the 128-dimensional cognitive manifold.

\section{SEC Vectors as Control Parameters}

SpiralBrain v3.0 uses four-dimensional SEC vectors as primary control signals:

\begin{itemize}
\item \textbf{Valence} ($v \in [-1,1]$): Controls positive/negative processing bias
\item \textbf{Arousal} ($a \in [0,1]$): Modulates processing intensity and speed
\item \textbf{Hazard Pressure} ($h \in [0,1]$): Regulates safety thresholds and risk assessment
\item \textbf{Neuromodulator Flexibility} ($f \in [-1,1]$): Adjusts pathway adaptability and rigidity
\end{itemize}

These vectors serve as foundational computational parameters, not descriptive labels. Emotional processing occurs at the core cognitive level, with SEC vectors directly modulating pathway activation and homeostasis controllers.

\section{Active Emotional Regulation System}

The active emotional regulator implements anticipatory intervention with four strategies:

\subsection{Coherence Prevention}
Threshold: 0.6, Effectiveness: 88\% \\
Maintains pathway alignment by preventing decoherence through targeted rebalancing.

\subsection{Elastic Preparation}
Threshold: 0.42, Effectiveness: 93\% \\
Prepares system for stress by adjusting communication decay and pathway damping.

\subsection{Meta-Priming}
Threshold: 0.62, Effectiveness: 75\% \\
Adapts safety thresholds and preloads homeostasis based on predicted needs.

\subsection{Cascade Prevention}
Threshold: 0.68, Effectiveness: 99\% \\
Prevents systemic failure through global damping and emergency pathway rebalancing.

Overall system effectiveness: 89\%, validating hypothesis H\textsubscript{7} (\geq75\% effectiveness threshold).

\section{Closed Cognitive Loop}

SpiralBrain implements a complete regulatory cycle:

\[
\text{Emotional State} \rightarrow \text{Prediction} \rightarrow \text{Intervention Plan} \rightarrow \text{Execution} \rightarrow \text{Learning} \rightarrow \text{Adaptation}
\]

This creates anticipatory regulation, preventing crises before they occur rather than responding after activation.

\section{Empirical Validation Metrics}

Core emotional foundation benchmark results demonstrate system maturity:

\begin{table}[H]
\centering
\caption{Emotional Regulation Performance Metrics}
\label{tab:emotional_metrics}
\begin{tabular}{@{}lcc@{}}
\toprule
Metric & Value & Interpretation \\
\midrule
Emotional Coherence & 0.431 & Active contextual transformation \\
Biological Accuracy & 1.000 & Perfect neurobiological constraint adherence \\
SEC Stability & 0.800 & High homeostasis enabling cognitive stability \\
Cognitive-Emotional Integration & 0.431 & Balanced coupling strength \\
Processing Performance & 42.9 seq/sec & Efficient native computation \\
\bottomrule
\end{tabular}
\end{table}

\section{Emotional Intelligence Quantification}

Emotional intelligence is measurable through SEC drift metrics:
- Optimal stability: SEC drift < 0.15
- Homeostatic maintenance: Consistent affective regulation across contexts
- Learning enhancement: 705\% improvement in emotional processing while preserving MMLU accuracy at 25.55\%

\section{Dynamical Systems Integration}

Emotions modulate cognitive trajectories within stable attractor basins:
- Low-positive Lyapunov exponents indicate controlled chaos
- SEC vectors define emotional state space boundaries
- Stability metrics quantify basin characteristics and recovery dynamics

\section{Implications for Cognitive Architecture}

Emotional regulation as control signals enables:
- Predictive homeostasis preservation
- Context-adaptive processing modulation
- Self-regulating cognitive stability
- Measurable emotional intelligence metrics

This approach demonstrates how emotions can serve as foundational computational substrates in synthetic cognition.

\section{Limitations}

Analysis constrained to SpiralBrain v3.0 instrumentation. No universality claims; results specific to observed SEC vector behavior. No proposals for alternative architectures; focus on demonstrated regulatory mechanisms.

All experimental results, analysis scripts, and logged outputs referenced in this study are publicly available in the SpiralBrain v3.0 repository \cite{SpiralBrainRepo}.

\section{Conclusion}

SpiralBrain v3.0 demonstrates emotions functioning as active control signals in synthetic cognition, achieving anticipatory regulation with 89\% effectiveness. SEC vectors serve as foundational computational parameters, enabling predictive homeostasis and adaptive processing modulation. This work establishes emotional regulation as a core architectural principle for stable synthetic cognition.

\section{Code Availability}

The experiments reported in this paper were conducted using the canonical SpiralBrain v3.0 configuration.
A public repository providing architectural documentation, configuration summaries, and reproducibility materials for this canonical configuration is available at: \\
\url{https://github.com/jhcragin/SpiralBrain-v3.0-public}. \\

The complete SpiralBrain codebase is not publicly available. Components beyond the public canonical configuration include research, experimental, and exploratory modules that were not exercised in the reported experiments and are maintained under a restricted research license.

\bibliographystyle{plain}
\bibliography{references}

\end{document}